(1) 波方程右端无时间delta函数，对界面附近积分得 \(u\) 与 \(u_t\) 连续；由于空间均匀，傅里叶波数 \(k\) 不变。界面之后频率模为 \(\pm\omega_2\)，其中 \(\omega_2=c_2k\)，不能要求频率在时间不均匀介质中守恒。
(2) 以复表示 \(u=a e^{ikx}[A e^{-i\omega_2t}+B e^{i\omega_2t}]\) 的实部，匹配 \(A+B=1\) 与 \(\omega_2(A-B)=\omega_1\)，故 \(A=(1+\omega_1/\omega_2)/2\)，\(B=(1-\omega_1/\omega_2)/2\)。
(3) 对实波的空间平均能量密度 \(\overline{\mathcal E}=\rho(u_t^2+c^2u_x^2)/2\)，界面前为 \(\rho a^2\omega_1^2/2\)，界面后为 \(\rho a^2\omega_2^2(A^2+B^2)/2\)。因此 \(\overline{\mathcal E}_+/\overline{\mathcal E}_-=[1+(c_2/c_1)^2]/2\)，差额由改变介质波速的外部驱动提供或吸收。
(4) \(c_2=c_1\) 时 \(A=1,B=0\) 且能量比为一。\(c_2\gg c_1\) 时两频率分支振幅接近 \(1/2\)，能量比增大；\(k\) 始终相同，体现时间界面保空间平移对称但破坏时间平移对称。
